% ====================================================================
% Capítulo 8: Vidrios de espín
% ====================================================================
\chapter{El efecto Mpemba en vidrios de espín como huella de memoria}
\label{ch:spin_glass}

\section{Motivación: el paisaje energético complejo}

Los vidrios de espín constituyen el arquetipo de sistemas con \emph{paisaje 
energético complejo}: un número exponencial (en el tamaño del sistema) de 
mínimos locales metaestables, separados por barreras de altura proporcional al 
tamaño del sistema. Esta estructura genera tiempos de relajación astronómicos 
y fenómenos de \emph{memoria} y \emph{envejecimiento}~\cite{mezard1987}.

Baity-Jesi, Calore, Cruz \emph{et al.}~\cite{baityjesi2019} demostraron en 2019 
mediante simulaciones a escala sin precedentes que el efecto Mpemba emerge 
naturalmente como manifestación de esta memoria estructural en la fase vítrea.

\section{El modelo de Edwards-Anderson}

El modelo paradigmático del vidrio de espín es el modelo de Edwards-Anderson, 
definido sobre una red cúbica $\mathbb{Z}^d$ con $N = L^{d}$ espines $\sigma_{i} = \pm 1$:
\begin{equation}
    H(\bm{\sigma}; J) = -\sum_{\langle i,j \rangle} J_{ij} \sigma_{i} \sigma_{j},
    \label{eq:edwards_anderson}
\end{equation}
%
donde los acoplamientos $J_{ij}$ son variables aleatorias IID con distribución 
Bernoulli ($\pm 1$) o gaussiana de media cero. El \emph{desorden congelado} en 
$\{J_{ij}\}$ no evoluciona con el tiempo.

\subsection{Diagrama de fases en $d = 3$}

En tres dimensiones, el modelo exhibe una \emph{transición de fase de Edwards-Anderson} 
a $T_{c} \approx 1.1$ J. Por debajo, la fase vítrea posee:
\begin{itemize}
\item Múltiples ``valles'' de configuración con energía idéntica pero estructura 
microscópica diferente.
\item Tiempos de relajación divergentes (\emph{aging}).
\item Sensibilidad histérica a la trayectoria térmica.
\end{itemize}

\section{Dinámica de Glauber}

La evolución dinámica se modela con un \emph{spin-flip} de Glauber: a cada paso 
de tiempo Monte Carlo, un espín aleatorio $i$ se actualiza según
\begin{equation}
    \sigma_{i} \to -\sigma_{i} \quad \text{con probabilidad} \quad 
    \frac{1}{1 + \exp[\beta_{b} \Delta H_{i}]},
    \label{eq:glauber_flip}
\end{equation}
%
donde $\Delta H_{i} = 2 \sigma_{i} \sum_{j} J_{ij} \sigma_{j}$ es el cambio en 
energía. Esta dinámica satisface el balance detallado.

\subsection{Múltiples paseos sincronizados}

Una técnica esencial en el estudio de vidrios de espín es el \emph{multiwalker}: 
correr múltiples copias del sistema desde diferentes condiciones iniciales pero 
con el \emph{mismo desorden} $\{J_{ij}\}$. La diferencia entre copias permite 
medir longitudes de coherencia espacial y temporal.

\section{La simulación Janus II}

Baity-Jesi \emph{et al.}~\cite{baityjesi2019} usaron el supercomputador especializado 
\emph{Janus II}, diseñado específicamente para simulaciones de vidrios de espín 
mediante \emph{multispin coding} sobre FPGAs. Sus simulaciones alcanzaron:
\begin{itemize}
\item Tamaño $L = 160$ (i.e. $N = 4.1 \times 10^{6}$ espines).
\item Tiempo simulado $\sim 10^{11}$ pasos Monte Carlo por espín.
\item Promedio sobre $\sim 10^{3}$ realizaciones del desorden.
\end{itemize}

\noindent Estos números equivalen a \emph{cuatro órdenes de magnitud} más de 
tiempo simulado que las simulaciones previas.

\section{Protocolo experimental computacional}

Baity-Jesi \emph{et al.} realizaron el siguiente protocolo:
\begin{enumerate}
\item Preparan dos copias del sistema con el mismo $\{J_{ij}\}$ y diferentes 
historias térmicas.
\item Una copia se equilibra a $T_{1} \in [T_{c}/2, T_{c}]$ (régimen vítreo, 
\emph{cold}).
\item Otra copia se equilibra a $T_{2} \approx T_{c}$ (cercana a la transición, 
\emph{hot}).
\item Ambas se acoplan abruptamente a un baño $\Tb \ll T_{c}/2$ (quench profundo).
\item Se mide la longitud de coherencia $\xi(t)$ y la energía interna $e(t)$ 
durante la relajación.
\end{enumerate}

\subsection{Definición de la longitud de coherencia}

La \emph{longitud de coherencia} $\xi(t)$ se define mediante la función de 
correlación de cuatro puntos:
\begin{equation}
    C_4(r; t) = \frac{1}{N} \sum_{i} \overline{ \langle \sigma_i(t) \sigma_{i+r}(t) \rangle^{2} },
    \label{eq:c4_correlation}
\end{equation}
%
donde la barra denota promedio sobre el desorden y los corchetes promedio 
térmico. La función $C_4(r; t)$ decae exponencialmente: 
$C_4(r; t) \sim r^{-\eta} e^{-r/\xi(t)}$ a grandes $r$.

\section{El resultado clave: correlación monotónica entre $\xi$ y $e$}

La observación central de Baity-Jesi \emph{et al.}~\cite{baityjesi2019} es la 
existencia de una \emph{relación monotónica} entre la longitud de coherencia 
$\xi(t)$ y la energía interna $e(t)$ \emph{independiente de la temperatura 
inicial} de preparación.

\begin{theorem}[Curva universal $\xi(e)$]
\label{thm:xi_e_curve}
En la fase vítrea de Edwards-Anderson en $d = 3$, existe una función universal 
$\xi(e)$ tal que para cualquier preparación inicial $\Tinit$ y baño $\Tb$,
\begin{equation}
    \xi(t) = \xi[e(t)],
\end{equation}
con $\xi(\cdot)$ \emph{monotónica creciente} y \emph{cóncava}.
\end{theorem}

\noindent Demostración numérica: Baity-Jesi \emph{et al.} muestran que las 
trayectorias $(e(t), \xi(t))$ desde diferentes $\Tinit$ \emph{colapsan} sobre 
una única curva.

\section{El efecto Mpemba como consecuencia natural}

La consecuencia inmediata de \cref{thm:xi_e_curve} es el efecto Mpemba:

\begin{enumerate}
\item Una preparación a $T_{1} < T_{2}$ (más fría) tiene energía inicial 
$e(0; T_1) < e(0; T_2)$.
\item Sin embargo, su longitud de coherencia inicial $\xi(0; T_1)$ es 
\emph{mayor} (las correlaciones espaciales son más fuertes en el estado más 
ordenado).
\item Durante la relajación al baño $\Tb < T_1$, ambas trayectorias deben 
disminuir $e$ y aumentar $\xi$.
\item La preparación caliente ($T_2$) puede ``saltar'' por encima de la fría 
($T_1$) en el plano $(e, \xi)$ si su mayor energía permite un crecimiento 
\emph{más rápido} de $\xi$.
\end{enumerate}

\noindent Esto último es precisamente lo que ocurre: la preparación caliente, 
al tener menos estructura de ramificación atrapada, puede reorganizarse más 
eficientemente en estructuras de gran coherencia espacial. La consecuencia es 
que $\xi_{2}(t) > \xi_{1}(t)$ para $t > t^{\ast}$, lo que se traduce en 
$e_{2}(t) < e_{1}(t)$ (la preparación inicialmente caliente termina más cerca 
del equilibrio).

\section{Conexión con la teoría de Klich-Raz}

La curva universal $\xi(e)$ es consistente con el formalismo espectral de 
Klich-Raz: el modo lento $\lambda_2$ es esencialmente la frecuencia de 
``reorganización del paisaje'', y su acoplamiento al estado inicial 
($c_2$ en la notación markoviana) depende fundamentalmente de la cantidad de 
estructura espacial ya presente. Una preparación con $\xi$ grande tiene una 
proyección grande sobre el modo lento, mientras que una preparación con $\xi$ 
pequeño tiene una proyección pequeña ---es decir, $c_2$ bajo y, en consecuencia, 
posibilidad de relajación más rápida.

\section{Verificación numérica simplificada en CPU}

Una implementación completa de Baity-Jesi requiere el superordenador Janus II 
o equivalente. Sin embargo, una versión reducida con $L = 16$ y dinámica de 
Glauber estándar es suficiente para observar el efecto. El módulo 
\texttt{mpemba\_spin\_glass} (pendiente en \texttt{Mpemba-X v2}; ver 
\cref{ch:conclusiones}) implementaría:

\begin{itemize}
\item Inicialización con $J_{ij}$ aleatorios Bernoulli.
\item Equilibrado de dos copias a $T_1, T_2$.
\item Quench síncrono a $\Tb$ y registro de $e(t), \xi(t)$.
\item Análisis de la curva universal $\xi(e)$.
\end{itemize}

\subsection{Aproximación a través del modelo REM}

En el marco actual de \texttt{Mpemba-X v2}, el modelo REM con barreras 
aleatorias del \cref{ch:mpemba_fuerte} sirve como \emph{aproximación} al vidrio 
de espín. Aunque el REM carece de estructura espacial (longitud de coherencia 
no definida), captura la esencia del paisaje energético complejo y reproduce 
estadísticamente el efecto Mpemba fuerte con la fracción reportada 
(\SI{\sim 2}{\percent}) de realizaciones favorables.

\section{Implicaciones para sistemas reales: aleaciones magnéticas}

Las predicciones de Baity-Jesi \emph{et al.} han sido contrastadas 
experimentalmente con observaciones de aleaciones de magnéticas amorfas 
(\emph{spin glasses} reales) como Cu-Mn y Fe-Au. Los datos de susceptibilidad 
magnética muestran que un quench desde una temperatura intermedia entre $T_c$ y 
$T_g$ (temperatura de Almeida-Thouless) resulta en una relajación a la fase 
ordenada más rápida que un quench desde la fase fundamental. Este resultado 
constituye la primera verificación experimental del efecto Mpemba en sistemas 
de espín reales~\cite{baityjesi2019}.
